```latex
\documentclass[11pt]{article}

\usepackage{amsmath,amssymb,amsthm,mathtools}
\usepackage[margin=1in]{geometry}
\usepackage{enumitem}
\usepackage[colorlinks=true,linkcolor=blue,citecolor=blue,urlcolor=blue]{hyperref}

\newtheorem{theorem}{Theorem}[section]
\newtheorem{proposition}[theorem]{Proposition}
\newtheorem{lemma}[theorem]{Lemma}
\newtheorem{corollary}[theorem]{Corollary}
\theoremstyle{definition}
\newtheorem{convention}[theorem]{Convention}
\newtheorem{definition}[theorem]{Definition}
\theoremstyle{remark}
\newtheorem{remark}[theorem]{Remark}

\newcommand{\Z}{\mathbb Z}
\newcommand{\Q}{\mathbb Q}
\newcommand{\F}{\mathbb F}
\newcommand{\Spin}{\operatorname{Spin}^{c}}
\newcommand{\HFhat}{\widehat{HF}}

\title{Prime summands in Dehn surgeries on \(L\)-space knots}
\author{}
\date{October 2, 2026}

\begin{document}
\maketitle

\begin{abstract}
With the standard convention that an \(L\)-space knot is a knot
\(K\subset S^{3}\) admitting a positive \(L\)-space surgery, the assertion that
every Dehn surgery on an \(L\)-space knot has at most two non-trivial prime
summands is not presently a theorem: it is precisely the \(L\)-space-knot case
of the two-summands problem for reducible surgery.  This note gives the
complete known reduction.  If a surgery on an \(L\)-space knot in \(S^{3}\)
has more than two prime summands, then the knot is hyperbolic, the slope is
\(2g(K)-1\), and the surgered manifold has the form
\[
        L(p,a)\# L(q,b)\# Y,
        \qquad pq=2g(K)-1,\quad (p,q)=1,\quad p,q>1,
\]
where \(Y\) is an integer homology sphere \(L\)-space.  Thus the original
question is equivalent to ruling out exactly this remaining case.  If
``\(L\)-space knot'' is interpreted more broadly as a knot in an \(L\)-space
homology sphere admitting an \(L\)-space surgery, then the answer is negative:
a regular fiber in the Poincare homology sphere admits a surgery to
\(L(2,1)\#L(3,2)\#L(5,4)\).
\end{abstract}

\section{Conventions and basic facts}

\begin{convention}
Unless explicitly stated otherwise, an \(L\)-space knot means a knot
\(K\subset S^{3}\) for which \(S^{3}_{n}(K)\) is an \(L\)-space for some
positive integer \(n\).  We write \(S^{3}_{r}(K)\) for \(r\)-surgery on \(K\).
Prime summands always means non-\(S^{3}\) prime summands in the
Kneser--Milnor decomposition.
\end{convention}

\begin{lemma}[Connected sums of \(L\)-spaces]
Let \(Y_{1},Y_{2}\) be rational homology spheres.  Then
\(Y_{1}\#Y_{2}\) is an \(L\)-space if and only if both \(Y_{1}\) and \(Y_{2}\)
are \(L\)-spaces.
\end{lemma}

\begin{proof}
Ozsvath and Szabo proved a Kunneth formula for Heegaard Floer homology:
over \(\F=\Z/2\Z\),
\[
 \HFhat(Y_{1}\#Y_{2},\mathfrak s_{1}\#\mathfrak s_{2})
 \cong
 \HFhat(Y_{1},\mathfrak s_{1})\otimes_{\F}
 \HFhat(Y_{2},\mathfrak s_{2});
\]
see \cite[Theorem 1.5]{OSClosed}.  Also
\[
 |H_{1}(Y_{1}\#Y_{2};\Z)|
 =
 |H_{1}(Y_{1};\Z)|\,|H_{1}(Y_{2};\Z)|.
\]
Thus, if both factors are \(L\)-spaces, the connected sum is an \(L\)-space.
Conversely, for every \(\mathfrak s\in\Spin(Y)\), \(\HFhat(Y,\mathfrak s)\)
is nonzero for a rational homology sphere \(Y\).  Hence all ranks in the
tensor product above are positive integers.  If the product rank is equal to
\(1\) for every pair \((\mathfrak s_{1},\mathfrak s_{2})\), then each factor
rank is equal to \(1\).  Therefore both \(Y_{1}\) and \(Y_{2}\) are
\(L\)-spaces.
\end{proof}

\section{The standard reducible-surgery input}

We use the following package of classical results.

\begin{theorem}[Reducible surgery package]\label{thm:package}
Let \(K\subset S^{3}\) be a nontrivial knot, and suppose \(S^{3}_{r}(K)\) is
reducible.  Then:
\begin{enumerate}[label=(\roman*)]
\item \(r\in \Z\) and \(r\neq 0\);
\item \(S^{3}_{r}(K)\) has a nontrivial lens-space summand;
\item \(S^{3}_{r}(K)\) has at most three nontrivial prime summands; if it has
three, then two are lens spaces and the third is an integer homology sphere.
\end{enumerate}
\end{theorem}

\begin{proof}
The integrality of a reducing slope is due to Gordon--Luecke
\cite{GLIntegral}.  The exclusion of \(r=0\) for nontrivial knots follows from
Gabai's Property \(R\) theorem \cite{Gabai}.  The existence of a nontrivial
lens-space summand is due to Gordon--Luecke \cite{GLReducible}.  The bound of
three summands, with the stated structure in the three-summand case, is
Howie's Corollary 5.3 in \cite{HowieSW}, combining work of Sayari and Valdez
Sanchez.
\end{proof}

\begin{lemma}\label{lem:homology}
Suppose
\[
        S^{3}_{n}(K)\cong L(p,a)\#L(q,b)\#Y,
\]
where \(Y\) is an integer homology sphere and \(p,q>1\).  Then
\[
        \gcd(p,q)=1,\qquad |n|=pq .
\]
\end{lemma}

\begin{proof}
For integer surgery on a knot in \(S^{3}\),
\[
        H_{1}(S^{3}_{n}(K);\Z)\cong \Z/|n|\Z .
\]
On the other hand,
\[
        H_{1}(L(p,a)\#L(q,b)\#Y;\Z)
        \cong \Z/p\Z\oplus \Z/q\Z .
\]
The direct sum \(\Z/p\Z\oplus \Z/q\Z\) is cyclic if and only if
\(\gcd(p,q)=1\), in which case it is cyclic of order \(pq\).  Hence
\(|n|=pq\).
\end{proof}

\section{Reduction for \(L\)-space knots in \(S^{3}\)}

\begin{lemma}[Non-hyperbolic \(L\)-space knots]\label{lem:nonhyp}
Let \(K\subset S^{3}\) be an \(L\)-space knot.  If \(K\) is not hyperbolic,
then every reducible surgery on \(K\) has at most two nontrivial prime
summands.
\end{lemma}

\begin{proof}
By Krcatovich, a knot in \(S^{3}\) admitting an \(L\)-space surgery is prime
\cite[Theorem 1.2]{Krcatovich}.  Thus a nontrivial non-hyperbolic \(L\)-space
knot is either a torus knot or a satellite knot.

If \(K\) is a torus knot, Moser's classification of surgeries on torus knots
shows that the unique reducible surgery is the cabling slope and gives a
connected sum of two lens spaces \cite{Moser}.

If \(K\) is a satellite knot, Scharlemann proved the cabling conjecture for
satellite knots \cite{Scharlemann}.  Thus \(K\) is a cable knot and the
reducing slope is the cabling slope.  If \(K=C_{u,v}(J)\) is the
\((u,v)\)-cable of \(J\), with \(u>1\), then the cabling-slope surgery has
the standard form
\[
        S^{3}_{uv}(C_{u,v}(J))
        \cong L(u,v)\# S^{3}_{v/u}(J).
\]
The slope \(v/u\) is nonintegral, so by Gordon--Luecke's integrality theorem
\cite{GLIntegral}, the second summand is irreducible unless it is \(S^{3}\).
Hence the surgery has at most two nontrivial prime summands.
\end{proof}

\begin{theorem}[Hom--Lidman--Zufelt]\label{thm:HLZ}
Let \(K\subset S^{3}\) be a hyperbolic \(L\)-space knot of genus \(g\).
If \(S^{3}_{r}(K)\) is reducible, then
\[
        r=2g-1 .
\]
\end{theorem}

\begin{proof}
This is \cite[Theorem 1.3]{HLZ}.
\end{proof}

\begin{theorem}[Ozsvath--Szabo]\label{thm:OSinterval}
Let \(K\subset S^{3}\) be a nontrivial \(L\)-space knot of genus \(g\).
Then every integer surgery with slope \(n\ge 2g-1\) is an \(L\)-space.
\end{theorem}

\begin{proof}
This is the integer-surgery part of the \(L\)-space surgery interval theorem
for knots in \(S^{3}\), proved by Ozsvath--Szabo in
\cite[Proposition 9.6]{OSLens}.
\end{proof}

\begin{theorem}[Precise reduction]\label{thm:reduction}
Let \(K\subset S^{3}\) be an \(L\)-space knot.  Suppose some surgery
\(S^{3}_{r}(K)\) has more than two nontrivial prime summands.  Then:
\begin{enumerate}[label=(\roman*)]
\item \(K\) is hyperbolic;
\item \(r=2g(K)-1\);
\item there exist relatively prime odd integers \(p,q>1\), lens spaces
\(L(p,a)\), \(L(q,b)\), and a nontrivial prime integer homology sphere
\(Y\), such that
\[
        S^{3}_{2g(K)-1}(K)\cong L(p,a)\#L(q,b)\#Y,
        \qquad pq=2g(K)-1;
\]
\item \(Y\) is an \(L\)-space.
\end{enumerate}
\end{theorem}

\begin{proof}
If \(r=\infty\), then \(S^{3}_{r}(K)\cong S^{3}\).  If \(r=0\), Gabai's
Property \(R\) theorem implies that reducibility forces \(K\) to be the unknot,
and \(0\)-surgery on the unknot is \(S^{1}\times S^{2}\), which is prime.
Thus \(r\neq 0,\infty\).

Assume \(S^{3}_{r}(K)\) has more than two nontrivial prime summands.  By
Theorem~\ref{thm:package}, \(r\) is a nonzero integer and the manifold has
exactly three nontrivial prime summands, two of which are lens spaces and the
third an integer homology sphere.  By Lemma~\ref{lem:nonhyp}, \(K\) cannot be
non-hyperbolic.  Hence \(K\) is hyperbolic.  Theorem~\ref{thm:HLZ} then gives
\[
        r=2g(K)-1.
\]

The structural statement from Theorem~\ref{thm:package} gives
\[
        S^{3}_{2g(K)-1}(K)\cong L(p,a)\#L(q,b)\#Y,
\]
where \(p,q>1\) and \(Y\) is an integer homology sphere.  Lemma~\ref{lem:homology}
implies
\[
        \gcd(p,q)=1,\qquad pq=2g(K)-1.
\]
Since \(2g(K)-1\) is odd, both \(p\) and \(q\) are odd.  The third summand
\(Y\) is nontrivial because the prime decomposition has three nontrivial
summands.

Finally, by Theorem~\ref{thm:OSinterval}, the surgery
\(S^{3}_{2g(K)-1}(K)\) is an \(L\)-space.  Lens spaces are \(L\)-spaces, and
Lemma 1.1 shows that in a connected sum of rational homology spheres, the
whole connected sum is an \(L\)-space if and only if every summand is an
\(L\)-space.  Therefore \(Y\) is an \(L\)-space.
\end{proof}

\begin{corollary}[Equivalence with the remaining open case]\label{cor:equiv}
For \(L\)-space knots in \(S^{3}\), the following are equivalent:
\begin{enumerate}[label=(\alph*)]
\item every Dehn surgery has at most two nontrivial prime summands;
\item there is no hyperbolic \(L\)-space knot \(K\) for which
\[
        S^{3}_{2g(K)-1}(K)
        \cong L(p,a)\#L(q,b)\#Y
\]
with \(p,q>1\), \(\gcd(p,q)=1\), \(pq=2g(K)-1\), and \(Y\) a nontrivial
integer homology sphere \(L\)-space.
\end{enumerate}
\end{corollary}

\begin{proof}
The implication \((a)\Rightarrow(b)\) is immediate.  Conversely, if \((a)\)
fails, then Theorem~\ref{thm:reduction} produces exactly a configuration of
the form excluded in \((b)\).
\end{proof}

\begin{remark}
Corollary~\ref{cor:equiv} is the precise present status of the original
question under the standard \(S^{3}\)-convention.  The two-summands
conjecture predicts that no reducible surgery on a knot in \(S^{3}\) has more
than two prime summands.  Theorem~\ref{thm:reduction} shows that the
\(L\)-space-knot case is reduced to a narrow hyperbolic case, but that case is
not presently ruled out in general.
\end{remark}

\section{Known Floer-theoretic obstruction in the remaining case}

The following theorem illustrates the strongest currently available
Floer-theoretic obstruction relevant to Theorem~\ref{thm:reduction}.

\begin{theorem}[Bodish--DeYeso]\label{thm:BD}
Let \(K\subset S^{3}\) be an \(L\)-space knot of genus \(g\).  Suppose
\[
        S^{3}_{2g-1}(K)\cong Y_{p}\#Y_{q},
\]
where \(p>q>1\), \(\gcd(p,q)=1\), and
\(|H_{1}(Y_{p};\Z)|=p\), \(|H_{1}(Y_{q};\Z)|=q\).  Then
\[
        q-1\le \operatorname{th}(K)\le \frac{p+q}{2},
\]
where \(\operatorname{th}(K)\) denotes the knot Floer thickness.  In particular,
if \(t\) is the smallest prime divisor of \(2g(K)-1\) and
\(\operatorname{th}(K)<t-1\), then \(K\) has no reducible surgery.
\end{theorem}

\begin{proof}
This is \cite[Theorem 1.9 and Corollary 1.11]{BodishDeYeso}.
\end{proof}

\section{A counterexample in the broader ambient-manifold sense}

If one allows \(L\)-space knots in arbitrary \(L\)-space homology spheres, the
answer to the question is negative.

\begin{theorem}\label{thm:counterexample}
There is a knot in an \(L\)-space integer homology sphere which admits an
\(L\)-space surgery with three prime summands.
\end{theorem}

\begin{proof}
Let \(P\) be the Poincare homology sphere.  It is a spherical manifold and is
an \(L\)-space.  Caudell proved that a regular fiber of the Seifert fibration
of \(P\) admits a Dehn surgery to
\[
        L(2,1)\#L(3,2)\#L(5,4);
\]
see \cite{Caudell}.  Each lens space is prime and is an \(L\)-space.  By
Lemma 1.1, their connected sum is an \(L\)-space.  Thus the regular fiber is
an \(L\)-space knot in the broader sense of a knot in an \(L\)-space homology
sphere admitting an \(L\)-space surgery, and the surgery has exactly three
nontrivial prime summands.
\end{proof}

\section*{Conclusion}

Under the standard convention \(K\subset S^{3}\), the requested affirmative
statement is open and is equivalent to excluding the hyperbolic configuration
in Corollary~\ref{cor:equiv}.  Under the broader convention allowing knots in
arbitrary \(L\)-space homology spheres, the statement is false by
Theorem~\ref{thm:counterexample}.

\begin{thebibliography}{99}

\bibitem{BodishDeYeso}
H. Bodish and R. DeYeso III,
\emph{Obstructing reducible surgeries: slice genus and thickness bounds},
arXiv:2209.01672.

\bibitem{Caudell}
J. Caudell,
\emph{Three lens space summands from the Poincare homology sphere},
arXiv:2101.01256.

\bibitem{Gabai}
D. Gabai,
\emph{Foliations and the topology of \(3\)-manifolds. III},
J. Differential Geom. \textbf{26} (1987), no. 3, 479--536.

\bibitem{GLIntegral}
C. McA. Gordon and J. Luecke,
\emph{Only integral Dehn surgeries can yield reducible manifolds},
Math. Proc. Cambridge Philos. Soc. \textbf{102} (1987), no. 1, 97--101.

\bibitem{GLReducible}
C. McA. Gordon and J. Luecke,
\emph{Reducible manifolds and Dehn surgery},
Topology \textbf{35} (1996), no. 2, 385--409.

\bibitem{HLZ}
J. Hom, T. Lidman, and N. Zufelt,
\emph{Reducible surgeries and Heegaard Floer homology},
Math. Res. Lett. \textbf{22} (2015), no. 3.

\bibitem{HowieSW}
J. Howie,
\emph{A proof of the Scott--Wiegold conjecture on free products of cyclic groups},
J. Pure Appl. Algebra \textbf{173} (2002), no. 2, 167--176.

\bibitem{Krcatovich}
P. Krcatovich,
\emph{The reduced knot Floer complex},
Topology Appl. \textbf{194} (2015), 171--201.

\bibitem{Milnor}
J. Milnor,
\emph{A unique decomposition theorem for \(3\)-manifolds},
Amer. J. Math. \textbf{84} (1962), 1--7.

\bibitem{Moser}
L. Moser,
\emph{Elementary surgery along a torus knot},
Pacific J. Math. \textbf{38} (1971), 737--745.

\bibitem{OSClosed}
P. Ozsvath and Z. Szabo,
\emph{Holomorphic disks and topological invariants for closed three-manifolds},
Ann. of Math. (2) \textbf{159} (2004), no. 3, 1027--1158.

\bibitem{OSLens}
P. Ozsvath and Z. Szabo,
\emph{On knot Floer homology and lens space surgeries},
Topology \textbf{44} (2005), no. 6, 1281--1300.

\bibitem{Scharlemann}
M. Scharlemann,
\emph{Producing reducible \(3\)-manifolds by surgery on a knot},
Topology \textbf{29} (1990), no. 4, 481--500.

\end{thebibliography}

\end{document}
```